% =============================================================================
%  Notas de clase — Sesión 20: Taller — Calculadora Financiera
%  Programación Avanzada — Universidad Panamericana
%  Compilar: pdflatex sesion20_taller_calculadora_financiera.tex
% =============================================================================
\documentclass[12pt, oneside]{article}
\usepackage{progavanzada}
\usepackage{array}
\usepackage{booktabs}

\begin{document}

\portada{Sesión 20}{Taller 3 — Calculadora financiera}{2026}

\tableofcontents
\newpage

% =============================================================================
\section{La calculadora más vendida de la historia}
% =============================================================================

En 1981, Hewlett-Packard lanzó la \textbf{HP-12C}: una calculadora del tamaño
de un bolsillo diseñada para profesionales de las finanzas.  Cuarenta y cinco
años después sigue fabricándose y vendiéndose, lo que la convierte en el
modelo de calculadora con mayor longevidad en la historia.  Hoy cuesta entre
70 y 100 dólares.  Corredores de bolsa, actuarios y analistas financieros la
llevan en el saco porque cumple una función muy concreta: calcular intereses,
flujos de efectivo y valores presentes sin perder tiempo arrancando una
laptop.

La HP-12C no tiene nada mágico.  Sus funciones más usadas son exactamente
las que implementarás en este taller: interés compuesto, valor del dinero en
el tiempo, y operaciones cotidianas como IVA y propinas.  La diferencia es
que la tuya correrá en la terminal y no costará 100 dólares.

\begin{tecnico}[La calculadora financiera en el mundo profesional]
  La \textbf{HP-12C} y la \textbf{Texas Instruments BA II Plus} (1991) son
  las dos calculadoras financieras dominantes en el mercado profesional.
  Ambas están permitidas en exámenes de certificación financiera como el
  \textit{CFA} (Chartered Financial Analyst) y el \textit{FRM} (Financial
  Risk Manager); otras calculadoras no están permitidas.  El hecho de que
  profesionales paguen decenas de dólares por hardware dedicado muestra
  cuánto valor tiene tener estas funciones disponibles al instante.
  \par\smallskip
  \textit{Referencia:} CFA Institute. (2024). \textit{CFA Program Approved
  Calculators}. \url{https://www.cfainstitute.org}
\end{tecnico}

% =============================================================================
\section{Objetivo del proyecto}
% =============================================================================

En este taller construirás una \textbf{calculadora financiera de consola} en
C++, organizada como un proyecto de múltiples archivos.  El programa mostrará
un menú con las funcionalidades disponibles, leerá los datos del usuario,
realizará los cálculos y mostrará el resultado con dos decimales de
precisión.

El proyecto aplica directamente lo de la sesión anterior: cada módulo
(\texttt{.h} + \texttt{.cpp}) tiene una responsabilidad delimitada, los
módulos se compilan por separado con un \texttt{Makefile}, y el ejecutable
se genera enlazando todos los objetos.

\begin{recuerda}
  Este proyecto no es una práctica de algoritmia ni de estructuras de datos.
  El objetivo principal es demostrar que sabes \textbf{organizar un programa
  en módulos} y \textbf{automatizar su compilación con Make}.  El contenido
  financiero es el vehículo, no el foco.
\end{recuerda}

% =============================================================================
\section{Arquitectura del proyecto}
% =============================================================================

El proyecto tiene \textbf{cuatro módulos mínimos}.  Puedes agregar más si
lo consideras necesario, pero no puedes reducirlos.

\begin{center}
\begin{tabular}{lp{9.5cm}}
  \toprule
  \textbf{Módulo} & \textbf{Responsabilidad} \\
  \midrule
  \texttt{main.cpp}
    & Punto de entrada.  Llama al menú principal en un bucle hasta
      que el usuario elija salir.  No contiene lógica de cálculo ni
      de presentación: solo orquesta los demás módulos. \\[4pt]
  \texttt{interfaz.h/.cpp}
    & Todo lo relacionado con la interacción en consola: mostrar el
      menú principal, mostrar los submenús de cada categoría, leer
      los datos que el usuario ingresa y mostrar los resultados
      formateados.  No realiza cálculos. \\[4pt]
  \texttt{matematicas.h/.cpp}
    & Funciones matemáticas de apoyo que utilizan los módulos
      financieros: potencia entera, cálculo de porcentaje y
      redondeo a dos decimales.  No contiene lógica financiera. \\[4pt]
  \texttt{financiero.h/.cpp}
    & Implementación de todas las funciones financieras del
      proyecto.  Utiliza las funciones de \texttt{matematicas.h}
      para los cálculos auxiliares.  No imprime nada en consola:
      recibe datos, devuelve resultados. \\
  \bottomrule
\end{tabular}
\end{center}

El grafo de dependencias entre módulos es:

\begin{verbatim}
  main.cpp
      |
  interfaz.h/.cpp
      |
  financiero.h/.cpp
      |
  matematicas.h/.cpp
\end{verbatim}

\begin{consejo}
  El módulo \texttt{interfaz} solo debe llamar funciones de \texttt{financiero}
  para obtener resultados y luego imprimirlos.  Los cálculos no van en la
  interfaz.  Esta separación hace que sea fácil cambiar la presentación
  (por ejemplo, pasar a interfaz gráfica) sin tocar la lógica.
\end{consejo}

\subsection{Qué debe contener cada módulo}

\subsubsection{main.cpp}

\begin{itemize}
  \item La función \texttt{main}.
  \item Un bucle que muestre el menú principal y procese la opción elegida
        hasta que el usuario indique que quiere salir.
  \item Inclusión de \texttt{interfaz.h}.
\end{itemize}

\subsubsection{interfaz.h / interfaz.cpp}

\begin{itemize}
  \item Una función que muestre el menú principal con todas las opciones
        disponibles numeradas.
  \item Una función por cada funcionalidad que lea los datos necesarios
        del usuario, llame a la función financiera correspondiente y
        muestre el resultado.
  \item Los resultados monetarios siempre con dos decimales y su unidad
        (e.g.\ ``\texttt{\$~1\,157.63 MXN}'').
\end{itemize}

\subsubsection{matematicas.h / matematicas.cpp}

\begin{itemize}
  \item Una función que calcule la potencia de un número real elevado a
        un exponente entero no negativo.  No uses \texttt{<cmath>} para
        implementarla: realiza el cálculo con un ciclo.
  \item Una función que calcule el porcentaje de un valor
        (e.g.\ \texttt{porcentaje(500, 16)} devuelve 80).
  \item Una función que redondee un número real a dos decimales.
\end{itemize}

\subsubsection{financiero.h / financiero.cpp}

\begin{itemize}
  \item Una función por cada una de las siete funcionalidades descritas
        en la sección siguiente.
  \item Cada función recibe los datos del cálculo y devuelve el resultado
        principal.  Si el cálculo produce más de un resultado (e.g.\ precio
        sin IVA y monto del IVA), usa parámetros de salida o una estructura.
  \item No imprime nada: toda la presentación queda en \texttt{interfaz}.
\end{itemize}

% =============================================================================
\section{Funcionalidades y fórmulas}
% =============================================================================

Las siguientes siete funcionalidades son \textbf{obligatorias}.

% ─── 4.1 ──────────────────────────────────────────────────────────────────────
\subsection{Funcionalidad 1 — Interés simple}

Cuando depositas dinero en una cuenta o pides un préstamo a \textbf{interés
simple}, el interés se calcula siempre sobre el capital original, sin importar
cuánto tiempo haya pasado.

\textbf{Variables:}
\begin{itemize}
  \item $C$ — capital inicial (monto depositado o prestado)
  \item $r$ — tasa de interés por período, expresada como decimal
        (e.g.\ 5\,\% → $r = 0.05$)
  \item $t$ — número de períodos (años, meses, etc.)
\end{itemize}

\textbf{Fórmulas:}
\[
  I = C \times r \times t
  \qquad
  M = C + I
\]

donde $I$ es el \textbf{interés generado} y $M$ es el \textbf{monto final}.

\begin{ejemplo}[Interés simple: \$1\,000 al 5\,\% anual, 3 años]
  \[
    I = 1{,}000 \times 0.05 \times 3 = 150
    \qquad
    M = 1{,}000 + 150 = 1{,}150
  \]
  Al cabo de 3 años habrás ganado \$150 y tendrás \$1\,150 en total.
\end{ejemplo}

% ─── 4.2 ──────────────────────────────────────────────────────────────────────
\subsection{Funcionalidad 2 — Interés compuesto}

Con \textbf{interés compuesto}, al final de cada período el interés se suma
al capital, y en el siguiente período el interés se calcula sobre el nuevo
total.  Es decir, ``ganas intereses sobre los intereses''.

\textbf{Variables:}
\begin{itemize}
  \item $C$ — capital inicial
  \item $r$ — tasa de interés por período (decimal)
  \item $n$ — número de períodos
\end{itemize}

\textbf{Fórmulas:}
\[
  M = C \times (1 + r)^n
  \qquad
  G = M - C
\]

donde $M$ es el \textbf{monto final} y $G$ es la \textbf{ganancia total}.

\begin{ejemplo}[Interés compuesto: \$1\,000 al 5\,\% anual, 3 años]
  \[
    M = 1{,}000 \times (1.05)^3 = 1{,}000 \times 1.157625 = 1{,}157.63
    \qquad
    G = 157.63
  \]
  Con interés simple habrías ganado \$150; con compuesto ganas \$157.63.
  La diferencia crece con el tiempo: a 30 años la brecha se vuelve enorme.
\end{ejemplo}

\begin{consejo}
  Para calcular $(1.05)^3$ usa la función de potencia de tu módulo
  \texttt{matematicas}.  No es necesario importar \texttt{<cmath>}
  para este proyecto.
\end{consejo}

% ─── 4.3 ──────────────────────────────────────────────────────────────────────
\subsection{Funcionalidad 3 — Calcular IVA}

El \textbf{Impuesto al Valor Agregado (IVA)} es un impuesto que se agrega al
precio de venta de casi todos los bienes y servicios.  En México la tasa
general es del \textbf{16\,\%}.

Esta funcionalidad recibe un precio sin IVA y calcula cuánto IVA se agrega
y cuál es el precio final.

\textbf{Variables:}
\begin{itemize}
  \item $p$ — precio sin IVA
  \item $t$ — tasa de IVA (decimal; en México $t = 0.16$)
\end{itemize}

\textbf{Fórmulas:}
\[
  \text{IVA} = p \times t
  \qquad
  p_{\text{total}} = p + \text{IVA} = p \times (1 + t)
\]

\begin{ejemplo}[IVA sobre \$500 a tasa del 16\,\%]
  \[
    \text{IVA} = 500 \times 0.16 = 80
    \qquad
    p_{\text{total}} = 500 + 80 = 580
  \]
\end{ejemplo}

% ─── 4.4 ──────────────────────────────────────────────────────────────────────
\subsection{Funcionalidad 4 — Desglosar IVA}

A veces un precio \textit{ya incluye} el IVA y necesitas saber cuánto
del total es IVA y cuánto es el precio base.  Esta es la operación inversa
a la funcionalidad anterior.

\textbf{Variables:}
\begin{itemize}
  \item $p_{\text{total}}$ — precio con IVA incluido
  \item $t$ — tasa de IVA (decimal)
\end{itemize}

\textbf{Fórmulas:}
\[
  p_{\text{sin IVA}} = \frac{p_{\text{total}}}{1 + t}
  \qquad
  \text{IVA} = p_{\text{total}} - p_{\text{sin IVA}}
\]

\begin{ejemplo}[Desglosar IVA de un precio de \$580 con tasa del 16\,\%]
  \[
    p_{\text{sin IVA}} = \frac{580}{1.16} = 500
    \qquad
    \text{IVA} = 580 - 500 = 80
  \]
\end{ejemplo}

\begin{recuerda}
  Las funcionalidades 3 y 4 son inversas: aplicar una y luego la otra
  debe devolverte el valor original.  Esto es útil para \textbf{verificar}
  tu implementación: si calculas el IVA sobre \$500 y obtienes \$580, y
  luego desglosas el IVA de \$580 y \textit{no} obtienes \$500, hay un error.
\end{recuerda}

% ─── 4.5 ──────────────────────────────────────────────────────────────────────
\subsection{Funcionalidad 5 — Cuenta de restaurante}

El programa calcula el total de una cuenta de restaurante, incluyendo la
propina, y lo divide entre los comensales.  Se asume que los \textbf{adultos}
pagan precio completo y los \textbf{niños} pagan \textbf{la mitad} de lo
que paga un adulto.

\textbf{Variables:}
\begin{itemize}
  \item $s$ — subtotal de la cuenta (suma de todos los platillos)
  \item $p$ — porcentaje de propina (e.g.\ 15 para 15\,\%)
  \item $n_a$ — número de adultos
  \item $n_n$ — número de niños
\end{itemize}

\textbf{Fórmulas:}
\[
  \text{propina} = s \times \frac{p}{100}
  \qquad
  T = s + \text{propina}
\]

Para dividir $T$ entre adultos y niños, donde cada niño paga la mitad de lo
que paga un adulto, sea $x$ el pago de un adulto:

\[
  n_a \cdot x + n_n \cdot \dfrac{x}{2} = T
  \qquad\Longrightarrow\qquad
  x = \dfrac{T}{n_a + \dfrac{n_n}{2}}
  \qquad
  x_n = \dfrac{x}{2}
\]

\begin{ejemplo}[Cuenta de \$480, propina 15\,\%, 2 adultos y 1 niño]
  \begin{align*}
    \text{propina} &= 480 \times 0.15 = 72\\
    T &= 480 + 72 = 552\\
    x &= \frac{552}{2 + 0.5} = \frac{552}{2.5} = 220.80 \text{ por adulto}\\
    x_n &= \frac{220.80}{2} = 110.40 \text{ por niño}\\
    \text{Verificación:} &\quad 2 \times 220.80 + 1 \times 110.40 = 552 \checkmark
  \end{align*}
\end{ejemplo}

% ─── 4.6 ──────────────────────────────────────────────────────────────────────
\subsection{Funcionalidad 6 — Descuento comercial}

Un \textbf{descuento comercial} reduce el precio de un artículo en un
porcentaje dado.

\textbf{Variables:}
\begin{itemize}
  \item $p_0$ — precio original
  \item $d$ — porcentaje de descuento (e.g.\ 25 para 25\,\%)
\end{itemize}

\textbf{Fórmulas:}
\[
  \text{ahorro} = p_0 \times \frac{d}{100}
  \qquad
  p_f = p_0 - \text{ahorro} = p_0 \times \left(1 - \frac{d}{100}\right)
\]

\begin{ejemplo}[Artículo de \$800 con 25\,\% de descuento]
  \[
    \text{ahorro} = 800 \times 0.25 = 200
    \qquad
    p_f = 800 - 200 = 600
  \]
\end{ejemplo}

% ─── 4.7 ──────────────────────────────────────────────────────────────────────
\subsection{Funcionalidad 7 — Conversión de divisas}

Convierte una cantidad de dinero de una moneda a otra, dado el tipo de
cambio.

\textbf{Variables:}
\begin{itemize}
  \item $A$ — cantidad en la moneda de origen
  \item $tc$ — tipo de cambio (cuántas unidades de la moneda destino
        equivalen a una unidad de la moneda origen)
\end{itemize}

\textbf{Fórmula:}
\[
  A_{\text{destino}} = A \times tc
\]

\begin{ejemplo}[100 USD a MXN con tipo de cambio 17.50]
  \[
    A_{\text{MXN}} = 100 \times 17.50 = 1{,}750.00 \text{ MXN}
  \]
\end{ejemplo}

\begin{consejo}
  El tipo de cambio que el usuario ingresa determina la dirección de la
  conversión.  Para convertir en sentido inverso (MXN → USD), el usuario
  debe ingresar el recíproco del tipo de cambio: si 1 USD = 17.50 MXN,
  entonces $tc = 1/17.50 \approx 0.0571$ para convertir de MXN a USD.
\end{consejo}

% =============================================================================
\section{Plan de pruebas}
% =============================================================================

Antes de hacer las capturas de pantalla para la entrega, verifica cada
funcionalidad con los valores de la siguiente tabla.  Si tu programa produce
los resultados esperados con estos datos, es muy probable que la lógica
sea correcta.

\begin{center}
\begin{tabular}{lp{5.5cm}p{4.5cm}}
  \toprule
  \textbf{Func.} & \textbf{Entradas} & \textbf{Resultado esperado} \\
  \midrule
  1 — Int.\ simple
    & $C = 1{,}000$, $r = 5$\,\%, $t = 3$
    & $I = 150.00$, $M = 1{,}150.00$ \\[4pt]
  2 — Int.\ compuesto
    & $C = 1{,}000$, $r = 5$\,\%, $n = 3$
    & $G = 157.63$, $M = 1{,}157.63$ \\[4pt]
  3 — Calcular IVA
    & $p = 500$, tasa $= 16$\,\%
    & IVA $= 80.00$, total $= 580.00$ \\[4pt]
  4 — Desglosar IVA
    & $p_{\text{total}} = 580$, tasa $= 16$\,\%
    & $p_{\text{sin IVA}} = 500.00$, IVA $= 80.00$ \\[4pt]
  5 — Restaurante
    & $s = 480$, propina $= 15$\,\%, 2 adultos, 1 niño
    & propina $= 72.00$, total $= 552.00$,
      adulto $= 220.80$, niño $= 110.40$ \\[4pt]
  6 — Descuento
    & $p_0 = 800$, $d = 25$\,\%
    & ahorro $= 200.00$, $p_f = 600.00$ \\[4pt]
  7 — Divisas
    & $A = 100$, $tc = 17.50$
    & $1{,}750.00$ \\
  \bottomrule
\end{tabular}
\end{center}

\begin{advertencia}
  Prueba también con entradas en los extremos: ¿qué pasa si el descuento
  es 0\,\%?  ¿Y si es 100\,\%?  ¿Qué pasa si no hay niños en la cuenta
  del restaurante ($n_n = 0$)?  Tu programa no debe crashear ni producir
  resultados sin sentido en estos casos.
\end{advertencia}

\begin{recuerda}
  Las funcionalidades 3 y 4 forman una prueba de consistencia: el resultado
  de desglosar IVA sobre el total que obtuviste en la funcionalidad 3 debe
  recuperar exactamente el precio original.
\end{recuerda}

% =============================================================================
\section{Entrega}
% =============================================================================

\subsection{Qué entregar}

Entregarás \textbf{un PDF} con capturas de pantalla de la terminal que
demuestren el funcionamiento del programa.  El código fuente no se entrega;
las capturas son la evidencia.

\subsection{Contenido del PDF}

El PDF debe incluir, en este orden:

\begin{enumerate}

  \item \textbf{Captura de compilación.}\\
        Muestra la ejecución de \texttt{make} desde cero
        (después de \texttt{make clean}) y la generación del ejecutable
        sin errores ni advertencias.

  \item \textbf{Una captura por funcionalidad} (siete en total).\\
        Cada captura debe mostrar:
        \begin{itemize}
          \item El menú con la opción seleccionada.
          \item Los datos que ingresaste.
          \item El resultado que imprimió el programa.
        \end{itemize}
        Para al menos \textbf{tres} funcionalidades, incluye
        \textbf{dos capturas distintas} (con datos diferentes) para
        mostrar que el programa funciona en casos distintos, no solo
        con los valores de la tabla de pruebas.

  \item \textbf{Captura del Makefile.}\\
        Muestra el contenido de tu \texttt{Makefile} en la terminal
        (\texttt{cat Makefile}).  Debe usar variables (\texttt{CXX},
        \texttt{CXXFLAGS}, \texttt{SRCS} u \texttt{OBJS}) y tener
        un objetivo \texttt{clean}.

\end{enumerate}

\subsection{Criterios de evaluación}

\begin{center}
\begin{tabular}{p{8cm}c}
  \toprule
  \textbf{Criterio} & \textbf{Puntos} \\
  \midrule
  El programa compila sin errores con \texttt{make}      & 10 \\
  El \texttt{Makefile} usa variables y tiene \texttt{clean} & 10 \\
  El proyecto tiene los cuatro módulos mínimos           & 10 \\
  Cada módulo respeta su responsabilidad (no mezcla lógica y presentación) & 10 \\
  \midrule
  Funcionalidades 1 y 2 (intereses) correctas            & 10 \\
  Funcionalidades 3 y 4 (IVA) correctas y consistentes   & 10 \\
  Funcionalidad 5 (restaurante) correcta                 & 15 \\
  Funcionalidades 6 y 7 (descuento y divisas) correctas  & 10 \\
  \midrule
  El PDF incluye todas las capturas requeridas           &  5 \\
  Al menos 3 funcionalidades con 2 casos de prueba distintos & 10 \\
  \midrule
  \textbf{Total}                                         & \textbf{100} \\
  \bottomrule
\end{tabular}
\end{center}

\begin{advertencia}
  Un programa que compila pero no tiene los módulos requeridos, o donde
  \texttt{financiero.cpp} imprime resultados directamente en lugar de
  devolverlos, pierde los puntos de arquitectura aunque las capturas
  muestren resultados correctos.
\end{advertencia}

% =============================================================================
\section*{Resumen}
% =============================================================================

\begin{itemize}
  \item El proyecto es una calculadora financiera de consola con \textbf{siete
        funcionalidades}: interés simple, interés compuesto, calcular IVA,
        desglosar IVA, cuenta de restaurante, descuento y conversión de divisas.
  \item El proyecto tiene \textbf{cuatro módulos}: \texttt{main},
        \texttt{interfaz}, \texttt{matematicas} y \texttt{financiero}.
  \item \texttt{financiero} calcula; \texttt{interfaz} presenta.
        Nunca al revés.
  \item El \texttt{Makefile} debe usar variables y compilar por etapas
        (\texttt{.cpp} → \texttt{.o} → ejecutable).
  \item La entrega es un \textbf{PDF con capturas} de cada funcionalidad
        y de la ejecución de \texttt{make}.
\end{itemize}

\end{document}
